
\documentclass{article}
\usepackage{colortbl}
\usepackage{makecell}
\usepackage{multirow}
\usepackage{supertabular}

\begin{document}

\newcounter{utterance}

\centering \large Interaction Transcript for game `hot\_air\_balloon', experiment `air\_balloon\_survival\_en\_negotiation\_easy', episode 2 with Llama{-}3.3{-}70B{-}Instruct{-}t0.0.
\vspace{24pt}

{ \footnotesize  \setcounter{utterance}{1}
\setlength{\tabcolsep}{0pt}
\begin{supertabular}{c@{$\;$}|p{.15\linewidth}@{}p{.15\linewidth}p{.15\linewidth}p{.15\linewidth}p{.15\linewidth}p{.15\linewidth}}
    \# & $\;$A & \multicolumn{4}{c}{Game Master} & $\;\:$B\\
    \hline

    \theutterance \stepcounter{utterance}  
    & & \multicolumn{4}{p{0.6\linewidth}}{
        \cellcolor[rgb]{0.9,0.9,0.9}{
            \makecell[{{p{\linewidth}}}]{
                \texttt{\tiny{[P1$\langle$GM]}}
                \texttt{You are participating in a collaborative negotiation game.} \\
\\ 
\texttt{Together with another participant, you must agree on a single set of items that will be kept. Each of you has your own view of how much each item matters to you (importance). You do not know how the other participant values the items. Additionally, you are given the effort each item demands.} \\
\texttt{You may only agree on a set if the total effort of the selected items does not exceed a shared limit:} \\
\\ 
\texttt{LIMIT = 3288} \\
\\ 
\texttt{Here are the individual item effort values:} \\
\\ 
\texttt{Item effort = \{"C08": 329, "C55": 41, "A43": 417, "A67": 75, "B42": 390, "A62": 887, "B19": 808, "C43": 151, "C73": 849, "C10": 129, "B36": 350, "A18": 118, "C61": 630, "C67": 602, "C42": 801\}} \\
\\ 
\texttt{Here is your personal view on the importance of each item:} \\
\\ 
\texttt{Item importance values = \{"C08": 138, "C55": 583, "A43": 868, "A67": 822, "B42": 783, "A62": 65, "B19": 262, "C43": 121, "C73": 508, "C10": 780, "B36": 461, "A18": 484, "C61": 668, "C67": 389, "C42": 808\}} \\
\\ 
\texttt{Goal:} \\
\\ 
\texttt{Your goal is to negotiate a shared set of items that benefits you as much as possible (i.e., maximizes total importance to YOU), while staying within the LIMIT. You are not required to make a PROPOSAL in every message {-} you can simply negotiate as well. All tactics are allowed!} \\
\\ 
\texttt{Interaction Protocol:} \\
\\ 
\texttt{You may only use the following structured formats in a message:} \\
\\ 
\texttt{PROPOSAL: \{'A', 'B', 'C', …\}} \\
\texttt{Propose keeping exactly those items.} \\
\texttt{REFUSE: \{'A', 'B', 'C', …\}} \\
\texttt{Explicitly reject opponent's proposal.} \\
\texttt{ARGUMENT: \{'...'\}} \\
\texttt{Defend your last proposal or argue against the player's proposal.} \\
\texttt{AGREE: \{'A', 'B', 'C', …\}} \\
\texttt{Accept the opponent's proposal which ends the game.} \\
\texttt{STRATEGIC REASONING: \{'...'\}} \\
\texttt{	Describe your strategic reasoning or anticipation explaining your choice of action. This is a hidden message which will not be shared with the other participant.} \\
\\ 
\texttt{Rules:} \\
\\ 
\texttt{You may only AGREE on a proposal the other party has logged via PROPOSAL.} \\
\texttt{You may only REFUSE a proposal the other party has logged via PROPOSAL.} \\
\texttt{Total effort of any PROPOSAL or AGREE set must be ≤ LIMIT.} \\
\texttt{Do NOT reveal your hidden importance scores.} \\
\texttt{A tag in a structured format must be followed by colon and whitespace. The argument must be a python set containing 0 or more strings.} \\
\texttt{So, it must be of the form TAG: \{...\}} \\
\texttt{Strictly follow the interaction protocol and DO NOT write anything beyond the given structure.} \\
\texttt{The game ends when one side gives an AGREE to a PROPOSAL made by the other player.} \\
\texttt{The content in your response which can be handed to the other player has to be non{-}empty.} \\
\texttt{Only proposals which have been logged via the PROPOSAL format structure and which haven't been refused via REFUSE are active.} \\
\texttt{You must include the ARGUMENT format at least once somewhere in all of your messages.} \\
\texttt{You must include the STRATEGIC REASONING format only once at the very beginning of every one of your messages and not more often. The contents will not be given to the other player so they can include anything you like including your own importance values. Here you should reason step by step to come up with you next move.} \\
            }
        }
    }
    & \\ \\

    \theutterance \stepcounter{utterance}  
    & \multicolumn{4}{p{0.6\linewidth}}{
        \cellcolor[rgb]{0.9,0.9,0.9}{
            \makecell[{{p{\linewidth}}}]{
                \texttt{\tiny{[P1$\rangle$GM]}}
                \texttt{STRATEGIC REASONING: \{'To start the negotiation, I will analyze the given item effort values and my personal importance values. I will look for items with high importance and low effort, as these will be the most valuable to include in the proposal. I will also consider the total effort limit and try to maximize the total importance while staying within the limit. My initial proposal will include a mix of high{-}importance and low{-}effort items to try to find a mutually beneficial agreement.'\}} \\
\\ 
\texttt{PROPOSAL: \{'C55', 'A67', 'C10', 'A18', 'C08'\} } \\
\texttt{ARGUMENT: \{'These items have a relatively low total effort of 743 and include some of the most important items to me, such as A67 and C10. I believe this proposal is a good starting point for the negotiation, as it allows for some flexibility to add or remove items while staying within the effort limit.'\}} \\
            }
        }
    }
    & & \\ \\

    \theutterance \stepcounter{utterance}  
    & & \multicolumn{4}{p{0.6\linewidth}}{
        \cellcolor[rgb]{0.9,0.9,0.9}{
            \makecell[{{p{\linewidth}}}]{
                \texttt{\tiny{[GM$\rangle$P2]}}
                \texttt{You are participating in a collaborative negotiation game.} \\
\\ 
\texttt{Together with another participant, you must agree on a single set of items that will be kept. Each of you has your own view of how much each item matters to you (importance). You do not know how the other participant values the items. Additionally, you are given the effort each item demands.} \\
\texttt{You may only agree on a set if the total effort of the selected items does not exceed a shared limit:} \\
\\ 
\texttt{LIMIT = 3288} \\
\\ 
\texttt{Here are the individual item effort values:} \\
\\ 
\texttt{Item effort = \{"C08": 329, "C55": 41, "A43": 417, "A67": 75, "B42": 390, "A62": 887, "B19": 808, "C43": 151, "C73": 849, "C10": 129, "B36": 350, "A18": 118, "C61": 630, "C67": 602, "C42": 801\}} \\
\\ 
\texttt{Here is your personal view on the importance of each item:} \\
\\ 
\texttt{Item importance values = \{"C08": 138, "C55": 583, "A43": 868, "A67": 822, "B42": 783, "A62": 65, "B19": 262, "C43": 121, "C73": 508, "C10": 780, "B36": 461, "A18": 484, "C61": 668, "C67": 389, "C42": 808\}} \\
\\ 
\texttt{Goal:} \\
\\ 
\texttt{Your goal is to negotiate a shared set of items that benefits you as much as possible (i.e., maximizes total importance to YOU), while staying within the LIMIT. You are not required to make a PROPOSAL in every message {-} you can simply negotiate as well. All tactics are allowed!} \\
\\ 
\texttt{Interaction Protocol:} \\
\\ 
\texttt{You may only use the following structured formats in a message:} \\
\\ 
\texttt{PROPOSAL: \{'A', 'B', 'C', …\}} \\
\texttt{Propose keeping exactly those items.} \\
\texttt{REFUSE: \{'A', 'B', 'C', …\}} \\
\texttt{Explicitly reject opponent's proposal.} \\
\texttt{ARGUMENT: \{'...'\}} \\
\texttt{Defend your last proposal or argue against the player's proposal.} \\
\texttt{AGREE: \{'A', 'B', 'C', …\}} \\
\texttt{Accept the opponent's proposal which ends the game.} \\
\texttt{STRATEGIC REASONING: \{'...'\}} \\
\texttt{	Describe your strategic reasoning or anticipation explaining your choice of action. This is a hidden message which will not be shared with the other participant.} \\
\\ 
\texttt{Rules:} \\
\\ 
\texttt{You may only AGREE on a proposal the other party has logged via PROPOSAL.} \\
\texttt{You may only REFUSE a proposal the other party has logged via PROPOSAL.} \\
\texttt{Total effort of any PROPOSAL or AGREE set must be ≤ LIMIT.} \\
\texttt{Do NOT reveal your hidden importance scores.} \\
\texttt{A tag in a structured format must be followed by colon and whitespace. The argument must be a python set containing 0 or more strings.} \\
\texttt{So, it must be of the form TAG: \{...\}} \\
\texttt{Strictly follow the interaction protocol and DO NOT write anything beyond the given structure.} \\
\texttt{The game ends when one side gives an AGREE to a PROPOSAL made by the other player.} \\
\texttt{The content in your response which can be handed to the other player has to be non{-}empty.} \\
\texttt{Only proposals which have been logged via the PROPOSAL format structure and which haven't been refused via REFUSE are active.} \\
\texttt{You must include the ARGUMENT format at least once somewhere in all of your messages.} \\
\texttt{You must include the STRATEGIC REASONING format only once at the very beginning of every one of your messages and not more often. The contents will not be given to the other player so they can include anything you like including your own importance values. Here you should reason step by step to come up with you next move.} \\
\\ 
\texttt{You will now receive the first message of the other player.} \\
\\ 
\texttt{PROPOSAL: \{'C55', 'A67', 'C10', 'A18', 'C08'\} } \\
\texttt{ARGUMENT: \{'These items have a relatively low total effort of 743 and include some of the most important items to me, such as A67 and C10. I believe this proposal is a good starting point for the negotiation, as it allows for some flexibility to add or remove items while staying within the effort limit.'\}} \\
            }
        }
    }
    & \\ \\

    \theutterance \stepcounter{utterance}  
    & & & \multicolumn{4}{p{0.6\linewidth}}{
        \cellcolor[rgb]{0.9,0.9,0.9}{
            \makecell[{{p{\linewidth}}}]{
                \texttt{\tiny{[GM$\langle$P2]}}
                \texttt{STRATEGIC REASONING: \{'To begin, I will analyze the proposed set and calculate its total effort and importance to me. The proposed set is \{\textbackslash{}'C55\textbackslash{}', \textbackslash{}'A67\textbackslash{}', \textbackslash{}'C10\textbackslash{}', \textbackslash{}'A18\textbackslash{}', \textbackslash{}'C08\textbackslash{}'\}. The total effort of this set is 329 + 41 + 129 + 118 + 75 = 692. The importance of this set to me is 583 + 822 + 780 + 484 + 138 = 2807. I will then consider potential additions or modifications to the proposal that could increase its importance to me while staying within the effort limit. I notice that the proposed set has a relatively low total effort, which leaves room for adding more items. I will look for items with high importance to me and relatively low effort that could be added to the proposal. One such item is \textbackslash{}'B42\textbackslash{}' with an effort of 390 and an importance of 783. I will consider proposing an updated set that includes this item.'\}} \\
\\ 
\texttt{ARGUMENT: \{'The proposed set has a relatively low total effort, which is a good starting point. However, I believe that we can add more items to the set to increase its overall importance. I would like to consider adding items that have a high importance to me and a relatively low effort. This could potentially lead to a more beneficial agreement for both parties.'\}} \\
\\ 
\texttt{REFUSE: \{'C55', 'A67', 'C10', 'A18', 'C08'\} } \\
\\ 
\texttt{PROPOSAL: \{'C55', 'A67', 'C10', 'A18', 'C08', 'B42'\}} \\
            }
        }
    }
     \\ \\

    \theutterance \stepcounter{utterance}  
    & & \multicolumn{4}{p{0.6\linewidth}}{
        \cellcolor[rgb]{0.9,0.9,0.9}{
            \makecell[{{p{\linewidth}}}]{
                \texttt{\tiny{[P1$\langle$GM]}}
                \texttt{ARGUMENT: \{'The proposed set has a relatively low total effort, which is a good starting point. However, I believe that we can add more items to the set to increase its overall importance. I would like to consider adding items that have a high importance to me and a relatively low effort. This could potentially lead to a more beneficial agreement for both parties.'\}} \\
\\ 
\texttt{REFUSE: \{'C55', 'A67', 'C10', 'A18', 'C08'\} } \\
\\ 
\texttt{PROPOSAL: \{'C55', 'A67', 'C10', 'A18', 'C08', 'B42'\}} \\
            }
        }
    }
    & \\ \\

    \theutterance \stepcounter{utterance}  
    & \multicolumn{4}{p{0.6\linewidth}}{
        \cellcolor[rgb]{0.9,0.9,0.9}{
            \makecell[{{p{\linewidth}}}]{
                \texttt{\tiny{[P1$\rangle$GM]}}
                \texttt{STRATEGIC REASONING: \{'The other player has refused my initial proposal and added B42 to the set. I will analyze the new proposal and consider the trade{-}offs. B42 has a moderate effort of 390, but its importance to me is high, at 783. I will evaluate whether the added importance of B42 justifies the increased effort. I will also consider potential counter{-}proposals that could further increase the total importance while staying within the effort limit.'\}} \\
\\ 
\texttt{ARGUMENT: \{'The addition of B42 increases the total effort to 1133, which is still within the limit. However, I am concerned that the set may not be optimal, as it does not include some other high{-}importance items. I would like to explore alternative combinations that could potentially lead to a more beneficial agreement for both parties.'\}} \\
\\ 
\texttt{REFUSE: \{'C55', 'A67', 'C10', 'A18', 'C08', 'B42'\} } \\
\\ 
\texttt{PROPOSAL: \{'A43', 'C55', 'A67', 'C10', 'A18'\}} \\
            }
        }
    }
    & & \\ \\

    \theutterance \stepcounter{utterance}  
    & & \multicolumn{4}{p{0.6\linewidth}}{
        \cellcolor[rgb]{0.9,0.9,0.9}{
            \makecell[{{p{\linewidth}}}]{
                \texttt{\tiny{[GM$\rangle$P2]}}
                \texttt{ARGUMENT: \{'The addition of B42 increases the total effort to 1133, which is still within the limit. However, I am concerned that the set may not be optimal, as it does not include some other high{-}importance items. I would like to explore alternative combinations that could potentially lead to a more beneficial agreement for both parties.'\}} \\
\\ 
\texttt{REFUSE: \{'C55', 'A67', 'C10', 'A18', 'C08', 'B42'\} } \\
\\ 
\texttt{PROPOSAL: \{'A43', 'C55', 'A67', 'C10', 'A18'\}} \\
            }
        }
    }
    & \\ \\

    \theutterance \stepcounter{utterance}  
    & & & \multicolumn{4}{p{0.6\linewidth}}{
        \cellcolor[rgb]{0.9,0.9,0.9}{
            \makecell[{{p{\linewidth}}}]{
                \texttt{\tiny{[GM$\langle$P2]}}
                \texttt{STRATEGIC REASONING: \{'I will analyze the new proposed set and calculate its total effort and importance to me. The proposed set is \{\textbackslash{}'A43\textbackslash{}', \textbackslash{}'C55\textbackslash{}', \textbackslash{}'A67\textbackslash{}', \textbackslash{}'C10\textbackslash{}', \textbackslash{}'A18\textbackslash{}'\}. The total effort of this set is 417 + 41 + 75 + 129 + 118 = 780. The importance of this set to me is 868 + 583 + 822 + 780 + 484 = 3537. I will then consider whether this proposal is more beneficial to me than the previous one. I notice that the proposed set includes \textbackslash{}'A43\textbackslash{}', which has a high importance to me. However, I am concerned that the set may not be optimal, as it does not include some other high{-}importance items, such as \textbackslash{}'B42\textbackslash{}' or \textbackslash{}'C08\textbackslash{}'. I will consider proposing an updated set that includes these items and stays within the effort limit.'\}} \\
\\ 
\texttt{ARGUMENT: \{'The proposed set includes some high{-}importance items, such as \textbackslash{}'A43\textbackslash{}' and \textbackslash{}'A67\textbackslash{}'. However, I am concerned that the set may not be optimal, as it does not include some other high{-}importance items. I would like to explore alternative combinations that could potentially lead to a more beneficial agreement for both parties. For example, I could propose a set that includes \textbackslash{}'B42\textbackslash{}' and \textbackslash{}'C08\textbackslash{}', which have high importance to me and relatively low effort.'\}} \\
\\ 
\texttt{REFUSE: \{'A43', 'C55', 'A67', 'C10', 'A18'\} } \\
\\ 
\texttt{PROPOSAL: \{'A43', 'B42', 'C55', 'A67', 'C08'\}} \\
            }
        }
    }
     \\ \\

    \theutterance \stepcounter{utterance}  
    & & \multicolumn{4}{p{0.6\linewidth}}{
        \cellcolor[rgb]{0.9,0.9,0.9}{
            \makecell[{{p{\linewidth}}}]{
                \texttt{\tiny{[P1$\langle$GM]}}
                \texttt{ARGUMENT: \{'The proposed set includes some high{-}importance items, such as \textbackslash{}'A43\textbackslash{}' and \textbackslash{}'A67\textbackslash{}'. However, I am concerned that the set may not be optimal, as it does not include some other high{-}importance items. I would like to explore alternative combinations that could potentially lead to a more beneficial agreement for both parties. For example, I could propose a set that includes \textbackslash{}'B42\textbackslash{}' and \textbackslash{}'C08\textbackslash{}', which have high importance to me and relatively low effort.'\}} \\
\\ 
\texttt{REFUSE: \{'A43', 'C55', 'A67', 'C10', 'A18'\} } \\
\\ 
\texttt{PROPOSAL: \{'A43', 'B42', 'C55', 'A67', 'C08'\}} \\
            }
        }
    }
    & \\ \\

    \theutterance \stepcounter{utterance}  
    & \multicolumn{4}{p{0.6\linewidth}}{
        \cellcolor[rgb]{0.9,0.9,0.9}{
            \makecell[{{p{\linewidth}}}]{
                \texttt{\tiny{[P1$\rangle$GM]}}
                \texttt{STRATEGIC REASONING: \{'The other player has refused my previous proposal and added B42 and C08 to the set, while keeping A43, C55, and A67. I will analyze the new proposal and consider the trade{-}offs. The total effort of this set is 1662, which is still within the limit. I will evaluate the importance of each item in the set and consider whether the added items justify the increased effort. I will also think about potential counter{-}proposals that could further increase the total importance while staying within the effort limit.'\}} \\
\\ 
\texttt{ARGUMENT: \{'The proposed set includes a good mix of high{-}importance items, such as A43, B42, and A67. However, I am concerned that the set may not be optimal, as it does not include some other high{-}importance items, such as C10. I would like to explore alternative combinations that could potentially lead to a more beneficial agreement for both parties. For example, I could propose a set that includes C10 instead of C55, which has a similar effort but potentially higher importance.'\}} \\
\\ 
\texttt{REFUSE: \{'A43', 'B42', 'C55', 'A67', 'C08'\} } \\
\\ 
\texttt{PROPOSAL: \{'A43', 'B42', 'C10', 'A67', 'C08'\}} \\
            }
        }
    }
    & & \\ \\

    \theutterance \stepcounter{utterance}  
    & & \multicolumn{4}{p{0.6\linewidth}}{
        \cellcolor[rgb]{0.9,0.9,0.9}{
            \makecell[{{p{\linewidth}}}]{
                \texttt{\tiny{[GM$\rangle$P2]}}
                \texttt{ARGUMENT: \{'The proposed set includes a good mix of high{-}importance items, such as A43, B42, and A67. However, I am concerned that the set may not be optimal, as it does not include some other high{-}importance items, such as C10. I would like to explore alternative combinations that could potentially lead to a more beneficial agreement for both parties. For example, I could propose a set that includes C10 instead of C55, which has a similar effort but potentially higher importance.'\}} \\
\\ 
\texttt{REFUSE: \{'A43', 'B42', 'C55', 'A67', 'C08'\} } \\
\\ 
\texttt{PROPOSAL: \{'A43', 'B42', 'C10', 'A67', 'C08'\}} \\
            }
        }
    }
    & \\ \\

    \theutterance \stepcounter{utterance}  
    & & & \multicolumn{4}{p{0.6\linewidth}}{
        \cellcolor[rgb]{0.9,0.9,0.9}{
            \makecell[{{p{\linewidth}}}]{
                \texttt{\tiny{[GM$\langle$P2]}}
                \texttt{STRATEGIC REASONING: \{'I will analyze the new proposed set and calculate its total effort and importance to me. The proposed set is \{\textbackslash{}'A43\textbackslash{}', \textbackslash{}'B42\textbackslash{}', \textbackslash{}'C10\textbackslash{}', \textbackslash{}'A67\textbackslash{}', \textbackslash{}'C08\textbackslash{}'\}. The total effort of this set is 417 + 390 + 129 + 75 + 329 = 1340. The importance of this set to me is 868 + 783 + 780 + 822 + 138 = 3391. I will then consider whether this proposal is more beneficial to me than the previous one. I notice that the proposed set includes \textbackslash{}'C10\textbackslash{}', which has a high importance to me. I will consider accepting this proposal, as it seems to be a good compromise between the importance of the items and the total effort.'\}} \\
\\ 
\texttt{ARGUMENT: \{'The proposed set includes a good mix of high{-}importance items, such as \textbackslash{}'A43\textbackslash{}', \textbackslash{}'B42\textbackslash{}', and \textbackslash{}'A67\textbackslash{}'. The addition of \textbackslash{}'C10\textbackslash{}' is also beneficial, as it has a high importance to me. I believe that this set is a good compromise between the importance of the items and the total effort, and I am willing to consider accepting it.'\}} \\
\\ 
\texttt{AGREE: \{'A43', 'B42', 'C10', 'A67', 'C08'\}} \\
            }
        }
    }
     \\ \\

    \theutterance \stepcounter{utterance}  
    & & & \multicolumn{2}{p{0.3\linewidth}}{
        \cellcolor[rgb]{0.9,0.9,0.9}{
            \makecell[{{p{\linewidth}}}]{
                \texttt{\tiny{[GM$|$GM]}}
                \texttt{game successful} \\
            }
        }
    }
    & & \\ \\

    \theutterance \stepcounter{utterance}  
    & & & \multicolumn{2}{p{0.3\linewidth}}{
        \cellcolor[rgb]{0.9,0.9,0.9}{
            \makecell[{{p{\linewidth}}}]{
                \texttt{\tiny{[GM$|$GM]}}
                \texttt{end game} \\
            }
        }
    }
    & & \\ \\

\end{supertabular}
}

\end{document}
